\section{The plane representation and its angular spectrum}\label{sec:angular}

The word kernels constructed in Section~\ref{sec:thermodynamic} now become
operators on a Hilbert space.  We identify their boundary values with local
spins and spatial translations, and their physical-angle coefficients with
the prescribed plane law.  The resulting spectral description will provide
both the normalization calculation and the separation into two chiral
limits.  Throughout this section
\[
 \pi/3<\lambda<\pi,\qquad \eta=\lambda/\pi,\qquad v=\pi/2,
 \qquad w=e^{iz}.
\]
All constants may depend on the fixed parameter.  The arguments below use
the ordinary homogeneous-row Hilbert norm.  No metric associated with a
diagonal cut, and no positive-association statement for $c<1$, is used.

\subsection{From word kernels to operators}

Let $\mathcal W$ be the complex vector space spanned by formal strict-row
words
\[
 A^{D_m}(z_m)\cdots A^{D_1}(z_1)\Omega,
 \qquad |\Re z_j|<v,
\]
including the empty word $\Omega$.  Give two words the limiting scalar
product obtained by reflecting the first word with the adjoint rule of
Lemma~\ref{lem:row-algebra}.  Proposition~\ref{prop:word-kernel} proves
existence of all these simultaneous limits.  The form is positive, since
every finite Gram matrix is a limit of positive finite-chain Gram matrices.
We quotient by its null space and complete; the resulting Hilbert space is
denoted by $\mathcal H$.

\begin{proposition}[The angular representation]\label{prop:angular-space}
On $\mathcal H$ the rows act as bounded holomorphic operator families, with
\[
 \|A^D(z)\|\le 2\|D\|,\qquad
 A^D(z)^*=A^{X\overline D X}(-\overline z).
\]
They satisfy the finite RTT exchange identities wherever their auxiliary
intertwiner is invertible.  There is a strongly continuous unitary group
$U_b$, $b\in\R$, fixing $\Omega$, such that
\begin{equation}\label{eq:angular-covariance}
 U_bA^D(z)U_b^*=A^D(z+ib).
\end{equation}
Words with real arguments in any nonempty open subinterval of $(-v,v)$
span a dense subspace.  The family $P=A^I$ consists of commuting normal
operators and has a simultaneous representation by almost surely
holomorphic scalar functions $f$ on the strip, with
\begin{equation}\label{eq:angular-scalar-star}
 f(-\overline z)=\overline{f(z)},\qquad |f(z)|\le2.
\end{equation}
In the half-plane coordinate this reality condition reads
$f(\overline w)=\overline{f(w)}$.
\end{proposition}

\begin{proof}
The uniform normalized row bounds of Lemma~\ref{lem:row-normalization},
applied before taking each Gram limit, show that left multiplication by a
row preserves null vectors and has the asserted bound.  Thus it extends
to the completion.  Adjoint and RTT identities pass to all word pairings
and then to bounded operators.  Matrix coefficients on words are
holomorphic by Proposition~\ref{prop:word-kernel}; local boundedness
extends this to all matrix coefficients and hence to operator
holomorphy.

Translate all arguments of a word by $ib$.  In its Gram product the
adjoint arguments are translated by the same amount, because
$-\overline{z+ib}=-\overline z+ib$.  Common-translation invariance of the
word kernels therefore makes this map an isometry.  Translation by $-ib$
is its inverse.  Continuity of the kernels gives strong continuity on
words, and the isometry bound extends it to $\mathcal H$.  This proves
\eqref{eq:angular-covariance}.  Notice that positive dilation of $w$
corresponds to one of the two signs of this unitary parameter.

If a vector is orthogonal to every word whose real arguments lie in a
fixed open interval, successive applications of the identity theorem to
the word arguments show that it is orthogonal to all strict words.  This
proves density.  Commutation of the untwisted rows and their adjoint rule
put them in one abelian von Neumann algebra.  Generate that algebra using
a countable dense set of arguments.  A multiplication representation,
together with the uniformly bounded Taylor series on a countable cover
of the strip, gives scalar analytic representatives on one common set
of full measure.  This also proves
\eqref{eq:angular-scalar-star}.  The spectral measure class is invariant
under imaginary translation by \eqref{eq:angular-covariance}.
\end{proof}

\subsection{Finite division and the two boundary lines}

Boundary rows have a direct spin interpretation.  Establishing that
interpretation in the limiting space requires more than convergence of
scalar eigenvalues.  We give the finite estimate and the common-radius
continuation argument explicitly.  This is the finite local mechanism
used in the angular-transfer construction of
\cite[Section ``Angular transfer and its two leading modes'']{AngularRepository2026};
the proof here checks it for the present parameter range.

\begin{lemma}[Boundary division and scalar continuation]
\label{lem:scalar-continuation}
There is $r_\lambda>0$, independent of the even chain length, for which
the two homogeneous quotients
\[
 T^D(z)T^I(z)^{-1},\qquad T^I(z)^{-1}T^D(z)
\]
are holomorphic and bounded by $C_\lambda\|D\|$ on discs of radius
$r_\lambda$ about $z=v$ and $z=-v$.  The limiting rows have strong
boundary values
\begin{equation}\label{eq:angular-endpoints}
 A^D(v)=SD_0,\qquad
 A^D(-v)=(XD^tX)_0S^{-1},
\end{equation}
where $S$ is the unitary site shift and $S\Omega=\Omega$.
Almost every spectral function extends analytically and without zeros
to bands of a common positive width about $\Re z=\pm v$, has modulus
one on these lines, and satisfies $|f|\le1$ between them.  Consequently
it is inner in the right half-plane, and every zero $\zeta$ satisfies
\begin{equation}\label{eq:angular-cone}
 \Re\zeta\ge c_0|\zeta|
\end{equation}
for a fixed $c_0>0$.  At each specified strict argument the spectral set
on which $f$ vanishes has measure zero.
\end{lemma}

\begin{proof}
Write $u=\eta(v-z)$, so $u=0$ at the upper endpoint, and put
$b(u)=\sin u/\sin\lambda$ and $r(u)=1-b(u)^2$.  Multiply the row
$T^D(u)$ by the unmarked row $\widetilde T^I(-u)$ with reversed
auxiliary order.  Folding the second auxiliary index by transposition
expresses the product as the two-auxiliary partial trace of
\[
 (D\otimes I)G_0(u)\cdots G_{L-1}(u),\qquad
 G_j(u)=R_{aj}(u)R_{a'j}(-u)^{t_{a'}}.
\]
The auxiliary line spanned by
$\Phi=2^{-1/2}(|++\rangle+|--\rangle)$ is invariant.  Inversion of
$R$ makes the action on this line equal to $r(u)I$.  At $u=0$ the
swap identity makes the range of $G_j(0)$ lie in
$\C\Phi\otimes\C^2$.  Thus, in the same auxiliary decomposition for
every site,
\begin{equation}\label{eq:folded-block}
 r(u)^{-1}G_j(u)=
 \begin{pmatrix}I&B_j(u)\\0&C_j(u)\end{pmatrix},\qquad
 \|B_j(u)\|\le C_\lambda,\quad
 \|C_j(u)\|\le C_\lambda|u|.
\end{equation}
These bounds follow directly from the entries of $R$, since
$\sin\lambda\ne0$.  For example, on a sufficiently small disc of
radius proportional to $\sin\lambda$ one has $|r|\ge3/4$ and
$\|C_j\|\le1/4$.  The upper-right block of the product in
\eqref{eq:folded-block} is bounded by a geometric series, while its
lower-right block has norm at most $4^{-L}$.  Taking the unmarked
auxiliary trace therefore gives
\[
 r(u)^{-L}T^I(u)\widetilde T^I(-u)=I+F_L(u),
 \qquad \|F_L(u)\|\le3\cdot4^{-L}.
\]
The off-diagonal triangular block has zero auxiliary trace.  The marked
partial trace is uniformly bounded by $C_\lambda\|D\|$.  The displayed
identity is invertible, so both finite square factors are invertible and
\[
 T^D(u)T^I(u)^{-1}
 =\bigl(r(u)^{-L}T^D(u)\widetilde T^I(-u)\bigr)(I+F_L(u))^{-1}
\]
has the required uniform bound.  Reversing the local product gives the
same argument with an invariant covector and proves the other quotient.
Crossing transfers it to the lower endpoint.  None of these steps uses
the sign of $\cos\lambda$.

The common finite vacuum is an eigenvector of $T^I$ also in these discs
by analytic continuation.  Its eigenvalue is nonzero there.  Applying
the quotient to that vacuum gives uniform analytic bounds on
$A_L^D(z)\Omega_L$, including all fixed Taylor orders.  Pairing with
strict words and taking normal-family limits yields a unique bounded
holomorphic vacuum vector on the disc, since its values on the interior
are already fixed by Proposition~\ref{prop:word-kernel}.

Here is a careful way to include endpoint vectors without assuming their
membership in $\mathcal H$.  First take a subsequential Gram enlargement
containing finite words of bounded endpoint rows and shifts.  Finite
swap and crossing identities give \eqref{eq:angular-endpoints} in this
enlargement.  Uniform Taylor estimates compare endpoint vacuum vectors
with nearby strict vectors in norm, before the chain limit as well as
after it.  Their limits are therefore the continued vectors just
constructed.

Choose a closed interval $I\subset(-v,v)$ with nonempty interior such
that every RTT intertwiner between either endpoint and an angle in
$I$ is invertible.  Such an interval exists: the excluded angles form a
finite set.  More explicitly, the possible zeros of the intertwiner
determinant occur at $s\in\pi\Z\pm\lambda$, where
$s=\eta(z-\alpha)$.  When $\lambda>\pi/2$ the possible interior
exceptions at the endpoints are
$\alpha=\pm(v-(\pi-\lambda)/\eta)$; avoiding these points suffices.
There is then a common endpoint disc on which all these inverses are
bounded.  Commute an endpoint-near row through a fixed word with angles
in $I$.  RTT expresses the result as a finite sum of bounded strict
rows acting on continued marked vacuum vectors.  The number and bounds
of the terms may depend on the word, but the radius of analyticity
does not.  The same finite-chain estimates give norm convergence to the
endpoint action.  Density of the $I$-words and the uniform interior row
bound extend this to strong convergence on all of $\mathcal H$.
In particular both $S$ and $S^{-1}$ preserve $\mathcal H$; their
restrictions are unitary inverses.  Successive strong approximations
give the same Gram limits for products containing several endpoint rows.

To pass from this vector continuation to scalar continuation with one
radius, choose countably many dense $I$-words $h_j$ and a finite measure
equivalent to the scalar spectral measure.  On a fixed endpoint disc
write the continued vector $P(z)h_j$ as
$\sum_{k\ge0}a_{jk}(z-v)^k$.  For $r'<R$ smaller than its common
radius, Cauchy's estimate gives
\[
 \|a_{jk}\|_{L^2}\le M_{j,R}R^{-k},\qquad
 \sum_{k\ge0}(r')^k\|a_{jk}\|_{L^2}<\infty.
\]
The series of fiber norms is therefore summable almost everywhere.
Intersecting over the countable words and radii produces pointwise
holomorphic vector series on the same disc.  On the interior part each
series equals $f(z)h_j$.  Its continuation is consequently proportional
to $h_j$.  At almost every spectral point at least one $h_j$ is nonzero,
and the resulting scalar ratios agree by the identity theorem.  This
constructs the scalar continuation on a common disc.

Equation~\eqref{eq:angular-endpoints} gives unitary values at both
endpoints.  Covariance and strong convergence give unitary boundary
values at $\pm v+ib$.  Applying this first on a countable dense set of
$b$'s and then using continuity shows that $|f|=1$ on each continued
boundary diameter.  Schwarz reflection gives
$f(z)\overline{f(2v-\overline z)}=1$ there, and the analogous formula
at $-v$.  Thus the continued functions are nonzero.  Rational imaginary
translations cover the boundary lines by discs of one radius, with a
common null set.  The bounded strip maximum principle now yields
$|f|\le1$.  Under $w=e^{iz}$ the excluded boundary bands give
\eqref{eq:angular-cone}.  Finally each scalar function has isolated
zeros.  Fubini on an imaginary-translation orbit gives a null zero set
at almost every point of that orbit, and invariance of the spectral
measure class transfers the assertion to every specified point.
\end{proof}

\subsection{Identification with the physical plane}

We next fix the probabilistic meaning of this representation.  The
physical row angle is $z=0$, or $w=1$, and the physical weights are
$cR(\lambda/2)$, with $a=b=1$.  Multiplying every vertex weight by $c$
does not affect any normalized finite-volume law.

\begin{proposition}[Physical law, both axes, and reflection]
\label{prop:physical-plane}
The row representation gives the local probabilities of the plane law
$\P_c$ in Theorem~\ref{thm:main}.  In particular its prescribed iterated
balanced-torus limit exists when $c<1$.  Quantization with time pointing
east and quantization with time pointing north describe this same law.
For an ordinary observable supported on one side of either coordinate
cut, the squared norm of its prepared state is its expectation paired
with its complex-conjugate reflected observable, with the corresponding
spin reflection.  Products of neutral height observables may be prepared
using paths confined to that side.
\end{proposition}

\begin{proof}
At fixed even transverse length, primitivity of the physical transfer
in the chosen sector identifies the long-torus limit with its normalized
Perron vector.  A finite collection of arrows is read by physical row
steps on a seam and by spins on finitely many equal-time sections.
Repeated seam marks are combined, and the spins are expressed by
\eqref{eq:angular-endpoints}.  Its joint probabilities are therefore
finite linear combinations of strict and endpoint word coefficients.
Their limits exist by Proposition~\ref{prop:word-kernel} and
Lemma~\ref{lem:scalar-continuation}.  When $c<1$ the root construction
used the exactly balanced sector, so these are precisely the two limits,
in their prescribed order, defining $\P_c$.

For completeness we identify the state also when the root construction
uses sectors tending to half filling.  Let $W_c$ be the physical vertex
weight and, for a stationary probability measure $\rho$ on allowed
arrow configurations, put
\begin{equation}\label{eq:angular-free-energy}
 \mathcal F_c(\rho)=-h(\rho)-\E_\rho\log W_c,
\end{equation}
where $h$ is specific entropy per vertex.  The equivalent gradient
formulation fixes the height parity and forgets an additive constant;
these changes have zero entropy cost per unit area.  Spatial translation
acts on gradients: after translating an anchored height representative,
one subtracts its new reference value.  Thus fixing the reference parity
does not remove any physical translation from the gradient action.
The domain and reference-height convention are those of
\cite[Section 2.3 and Appendix B]{Sheffield2005}; the entropy--energy
decomposition is also explicit in
\cite[Sections 2.3.2--2.3.4]{LammersTassy2024}.
Since the alphabet
is finite and $0<c<\infty$, the entropy and local energy in
\eqref{eq:angular-free-energy} are finite.  Hard constraints are included
by restricting to allowed configurations.

Write $U(0)$ for the homogeneous harmonic normalization in
Lemma~\ref{lem:row-normalization}, and set
$p_*=U(0)+\log c$.  The same boundary-unitarity and strip argument used
there applies to a row with specified auxiliary input and output spins:
it is a matrix-unit auxiliary mark.  A rectangle of width $a$ and
height $b$, with any boundary data, therefore has partition sum at most
\[
 \exp\{abp_*+C(a+b)\}.
\]
Indeed each open row has operator norm at most $2e^{ap_*}$, and
summing the boundary spins costs only an exponential in its perimeter.
The finite entropy variational inequality on the rectangle then gives
$h(\rho)+\E_\rho\log W_c\le p_*$ for every stationary allowed state.

For the near-balanced tori, choose their longitudinal lengths sufficiently
large to realize the Perron asymptotic.  Their constrained log partition
functions per vertex tend to $p_*$ by
Lemma~\ref{lem:row-normalization}.  Tile each torus by fixed-size squares.
Entropy subadditivity, with a negligible uncovered fraction, bounds its
entropy per vertex above by the entropy of a fixed-square marginal
divided by the square area, plus a vanishing boundary error.  Pass to the
local limit at this fixed square size, then let the square size grow.
The exact finite identity ``entropy plus mean log weight equals log
partition function'' proves the reverse inequality
$h(\rho)+\E_\rho\log W_c\ge p_*$.  Hence the limiting state attains
the global minimum of \eqref{eq:angular-free-energy}.  Its expected
spatial tilt is zero by the sector sizes.  The other expected tilt is
zero as well: at the physical angle the current mark is skew-adjoint,
whereas its finite physical expectation is real.  This observation
specifies the tilt; it is not being used as a uniqueness argument.

The equilibrium characterization following equation (6) in
\cite{DKLM2026}, for $1\le c\le2$, identifies the
finite-specific-free-energy minimizer with the balanced plane law.
We spell out the variational-domain reduction needed to apply this
characterization, since a zero expected tilt alone would not suffice.
For a fixed spatial translation lattice $\mathcal L$, let
$\mathcal P_{\mathcal L}$ denote invariant laws of allowed gradients and
write
\[
 \sigma_{\mathcal L}(a)
  =\inf\{\mathcal F_c(\rho):\rho\in\mathcal P_{\mathcal L},\ S(\rho)=a\}.
\]
Here $S(\rho)$ is the expected physical slope.  Both the preceding
perimeter bound and the tiling proof apply on this domain, using the
average energy over a fundamental cell if necessary.  We keep the same
$\mathcal L$ when forming the infimum and the ergodic decomposition.
The surface tension is unchanged on passing from full translations to
a finite-index sublattice: average any sublattice-invariant law over
the finitely many translation cosets.  This preserves its slope and
its entropy--energy density and makes it fully invariant; the reverse
inequality between the infima follows from inclusion of the domains.
Thus one may use the full physical gradient translation action.  If a
periodic height-potential representation uses a parity-preserving
sublattice, the argument is performed throughout on that sublattice
and is then projected back to arrows.

Specific free energy is affine under this ergodic decomposition
\cite[Chapter 3]{Sheffield2005}.  Thus every component of a globally
minimizing law is itself a global minimizer: each component has free
energy at least the global infimum, and their average equals it.
For $c\ge1$, strict convexity of the six-vertex surface tension on the
interior of its slope domain is
\cite[Theorem 13.16 of the 2020 preprint]{LammersTassy2024}.
Its stated condition is $c^2\ge\max(1,1)$, and the accompanying
simply-attractive-potential description supplies the height-model
setting.  Arrow reversal makes $\sigma_{\mathcal L}$ even.  The origin
is an interior slope: height functions formed as sums of two
one-dimensional unit-increment paths give allowed stationary gradient
laws at every slope in $[-1,1]^2$, all of finite specific free energy.
Hence strict convexity implies
$\sigma_{\mathcal L}(a)>\sigma_{\mathcal L}(0)$ for every nonzero
interior $a$.  A boundary slope cannot be another minimizer: otherwise
convexity along the segment from $0$ to that slope would contradict
this strict inequality at an interior point.  Every minimizing ergodic
component therefore has slope zero and is a fixed-slope minimizer.
The variational principle makes it a minimal gradient phase.

At this last step we use the zero-slope uniqueness supplied by the
combined equilibrium characterization in \cite{DKLM2026}, with its
Sheffield and delocalization inputs and its range $1\le c\le2$.
This is not an assertion of uniqueness for every Gibbs state, nor an
extension to $c>2$.  The domain reduction above explains why the
possibly nonergodic auxiliary state is covered, even when a formulation
of the underlying height result is stated for minimal phases.
The gradient-to-arrow correspondence is one-to-one after forgetting
the additive height constant, so their common law is $\P_c$.

It remains to compare the two axes in the range $c<1$, where we have not
used an equilibrium uniqueness theorem.  Fix a rectangle.  Orient its
input path upward along the left side and then eastward along the top
side, and express all spins relative to those directed steps.  Let $K$
be the tensor of its unmarked physical tiles, viewed as a matrix from
the input sides to the other two sides.  It is a product of physical
$R(\lambda/2)$ matrices and permutations of tensor factors.  The local
eigenvalues before the harmless factor $c$ are
$1/c,1/c,1/c+1,1/c-1$.  Thus $K$ is invertible for $c<1$.

For a finite surrounding cylinder, cut on all four sides and contract
the exterior with its normalized vacuum data.  This gives a linear
functional $L$ of the inside matrix, with $L(K)=1$.  The normalization
includes the exterior transfer eigenvalues appropriate to that cylinder.
Putting elementary matrices $D_1,\ldots,D_q$ on the input path tests
the functional on $K(D_1\otimes\cdots\otimes D_q)$: the row index of
each mark is inside and its column index outside.  Since tensor products
of elementary matrices span the input matrix algebra and $K$ is fixed
and invertible, the limiting marked path values determine every entry
of the limiting exterior functional.  They also prove that those
entries have limits.

In the eastward representation the successive path angles are $v,0$;
in the northward representation, with space pointing west, they are
$0,-v$.  To compare them, first replace the first pair by
$(v-\varepsilon,\varepsilon)$ and translate both angles by $-v$.
The resulting pair is $(-\varepsilon,-v+\varepsilon)$, still strict.
Common-translation invariance identifies the marked word values.
Strong boundary convergence then permits $\varepsilon\downarrow0$.
The spin and matrix-index conventions agree: at a reversed collapsed
step the mark is $XD^tX$, exactly as in
\eqref{eq:angular-endpoints}.  Crossing only relabels those spins; the
interior physical tensor remains $K$.  Thus the two normalized limiting
functionals agree on $KD$ for every elementary tensor $D$, and hence on
every inside matrix.  All local arrow probabilities agree.  This
argument compares normalized exterior functionals; it requires no
equality between eigenvalues of cylinders of different circumferences.

Finally, in either homogeneous axial quantization, transposing the
finite row product at a coordinate cut reflects its tiles and changes
the spin mark by the stated adjoint rule.  The finite Gram square is
therefore the expectation of the observable and its conjugate reflected
copy.  Pass to the local and Gram limits to obtain the assertion.
Neutral height differences can be read along paths on the prescribed
side of the cut, so the same identity applies to their finite products.
\end{proof}

\subsection{Finite inner functions and a measure on dilation orbits}

We now restrict when indicated to the closed diagonal-word subspace
$\mathcal D\subset\mathcal H$: it is generated by words all of whose
marks are diagonal.  It is invariant under $P,J$ and their adjoints,
where $J=A^Z$ and $Z=\diag(1,-1)$.  Endpoint diagonal rows show that it
contains the states needed for ordinary height insertions.  Write
$P(w)$ and $J(w)$ in half-plane coordinates from now on.

\begin{lemma}[Inner spectra and the current measure]\label{lem:inner-spectra}
On $\mathcal D$, almost every spectral function is either $1$ or a
finite Blaschke product in the right half-plane, with
\begin{equation}\label{eq:finite-inner-symmetry}
 f(0)=f(\infty)=1,\qquad
 f(\overline w)=\overline{f(w)}.
\end{equation}
Its zeros obey \eqref{eq:angular-cone}.  In particular its degree is
even; write it as $2n_f$ when it is nonconstant.  Every zero produces
a fixed-width, fixed-depth dip of $|f(e^s)|$ near its logarithmic radius.

There is a possibly infinite positive measure $\mu$ on the nonconstant
functions and a measurable spectral vector $\mathsf h$ such that
\begin{equation}\label{eq:current-cocycle}
 J(w)\Omega=(1-f(w))\mathsf h,
 \qquad \mu(\dd f)=\|\mathsf h(f)\|^2\,\sigma(\dd f).
\end{equation}
Here $\sigma$ is any scalar measure used for the direct-integral
representation; the resulting measure $\mu$ does not depend on that
choice.  It is invariant under dilation of the argument.  On the
cross-section where the smallest zero radius is one, it has the form
\begin{equation}\label{eq:orbit-disintegration}
 \mu=\dd s\,\nu(\dd F),\qquad f(w)=F(e^s w),\qquad
 \kappa:=\nu(1)<\infty.
\end{equation}
If $\mathcal L(F)$ is the difference between its largest and smallest
logarithmic zero radii, then
\begin{equation}\label{eq:orbit-width-tail}
 \nu\{\mathcal L>t\}\le Ce^{-2\eta t}\quad(t\ge1).
\end{equation}
For a fixed diagonal word $h$,
\begin{equation}\label{eq:diagonal-boost-tail}
 \|(1-P(e^s))h\|\le C_h e^{-\eta|s|}.
\end{equation}
For products of fixed words placed at scales $r_i$, the corresponding
bound outside their fixed internal argument ranges is
$C\sum_i e^{-\eta|s-\log r_i|}$, uniformly when these ranges separate.
\end{lemma}

\begin{proof}
As an auxiliary argument is dilated far from all arguments in a fixed
word, the scalar-normalized RTT intertwiner tends to an invertible
diagonal matrix, with error $O(e^{-\eta|s|})$.  Its limiting diagonal
matrix commutes with the diagonal joint auxiliary marks.  Move $P(e^s)$
through the word to the vacuum, where it equals the identity.  The
exchange errors are finite sums of bounded marked-row products, proving
\eqref{eq:diagonal-boost-tail}.  Moving through several separated groups
gives the stated sum, since the number of factors is fixed.  These are
operator-norm bounds on each finite exchange error, not just bounds on
its vacuum expectation.

On a countable total set in $\mathcal D$, the spectral integral of
$|1-f(e^s)|^2$ over $s$ is consequently finite.  The general inner
factorization allowed by Lemma~\ref{lem:scalar-continuation} is
\begin{equation}\label{eq:general-inner-factorization}
 f(w)=\gamma e^{-aw-b/w}\prod_j\beta_j(w),\qquad
 a,b\ge0,\quad |\gamma|=1,
\end{equation}
where $\beta_j$ is a constant-phase multiple of
$(w-\zeta_j)/(w+\overline\zeta_j)$ and
$\sum_j\Re\zeta_j/(1+|\zeta_j|^2)<\infty$.
Indeed continuation through every nonzero finite point of the imaginary
axis excludes singular measure there; the only remaining singular
atoms are at $0$ and $\infty$, giving $b/w$ and $aw$.
If $a>0$ or $b>0$, then $|f(e^s)|$ tends to zero at one end, contrary
to the integrability just proved.  For a zero $\zeta$ with
$r=|\zeta|$, the identity
\[
 \left|\frac{t-\zeta}{t+\overline\zeta}\right|^2
 =1-\frac{4t\Re\zeta}{|t+\overline\zeta|^2},\qquad t>0,
\]
and \eqref{eq:angular-cone} imply
$1-|f(e^s)|\ge\varepsilon_0$ whenever
$|s-\log r|\le\delta_0$, with fixed positive
$\varepsilon_0,\delta_0$.  Infinitely many zeros would escape to one
end and supply infinitely many disjoint such intervals, again
contradicting integrability.  The resulting finite product has limits
at both ends, and integrability forces both limits to be one.
Reality was already established.  Nonreal zeros occur in conjugate
pairs, and the endpoint quotient is $(-1)^{\deg f}$, so the degree is
even.  This proves \eqref{eq:finite-inner-symmetry} and the dip assertion.

Differentiate commutation of equal diagonal-twist transfers at twist
zero and apply it to $\Omega$.  Since $P(w)\Omega=\Omega$, it gives
\begin{equation}\label{eq:twist-cocycle}
 (1-P(w'))J(w)\Omega=(1-P(w))J(w')\Omega.
\end{equation}
There is no component of $J(w)\Omega$ on $f=1$.  Indeed, let $Q$ be
the projection to that spectral set, which contains $\Omega$.
The endpoint identities and their dilates glue $QJ(w)Q$ to the left
half-plane by $QJ(-w)Q=-QJ(w)Q$, because $P(w)=I$ on $Q$.
It is bounded on the punctured plane, so the singularities at zero
and infinity are removable.  Liouville's theorem and oddness make it
zero.  On every
nonconstant fiber $|f(1)|<1$, so one may set
$\mathsf h=J(1)\Omega/(1-f(1))$.  Equation~\eqref{eq:twist-cocycle}
gives \eqref{eq:current-cocycle}, first at a countable set of arguments
and then everywhere, including endpoints by strong convergence.
Covariance, with $U_b\Omega=\Omega$, makes this quotient transform
without a scalar multiplier.  Its squared fiber norm therefore gives
a dilation-invariant measure.

Every nonconstant finite product has a unique smallest zero radius.
Rescale it to one to obtain a measurable cross-section.  Translation
invariance in its logarithmic dilation coordinate gives
\eqref{eq:orbit-disintegration}, initially as a sigma-finite
disintegration.  Each orbit has a dip of fixed width at its first zero,
where $|1-f(1)|\ge\varepsilon_0$.  Since
\[
 \int |1-f(1)|^2\dd\mu=\|J(1)\Omega\|^2<\infty,
\]
this disintegration is locally finite and $\kappa<\infty$.

Applying \eqref{eq:diagonal-boost-tail} to $J(1)\Omega$ gives
\begin{equation}\label{eq:two-dip-bound}
 I(l):=\int |1-f(1)|^2|1-f(e^l)|^2\dd\mu
 \le Ce^{-2\eta l}\qquad(l\ge1).
\end{equation}
On an orbit of width $\mathcal L$, integrate first over the dip window
of its smallest zero and then over $l$ placing the second argument in
the dip window of its largest zero.  Both intervals have a fixed width
and both factors a fixed lower bound.  Thus
\[
 c\,\nu\{\mathcal L>t\}
 \le \int_{t-2\delta_0}^{\infty} I(l)\dd l,
\]
after changing the constant for bounded $t$.  This proves
\eqref{eq:orbit-width-tail}.
\end{proof}

\subsection{Low spectral bands and the two leading current modes}

The next estimate is an operator statement on the full space
$\mathcal H$.  It therefore applies also when the states on either side
of a current are different.  Define
\[
 H_0=-\log|P(1)|,\qquad Q_u=\bbone_{[0,u]}(H_0),\qquad
 Q_u^\sigma=Q_u\bbone_{\{\sgn P(1)=\sigma\}},\quad \sigma=\pm1.
\]
The logarithm is defined by spectral calculus; the fixed-argument
no-zero assertion makes its value finite almost everywhere.  Put
\begin{equation}\label{eq:angular-N-definition}
 N(w)=\tfrac12\{J(w)P(w)+P(w)J(w)\}.
\end{equation}
This current averages the two positions of a mark among two consecutive
steps.  The factor $1/2$ will be retained throughout.

\begin{lemma}[Low-band current estimates]\label{lem:low-band}
There are constants $A,a,C,u_0>0$ such that, for $0<u<u_0$,
$Q_uJ(w)Q_u$ extends to a bounded holomorphic function on the full
annulus $Au<|w|<a/u$.  If
\begin{equation}\label{eq:angular-Laurent}
 Q_uJ(w)Q_u=\sum_{j\in\Z}w^j J_{u,j},
\end{equation}
then $\|J_{u,j}\|\le C(Cu)^{|j|}$, with the exponent zero for $j=0$.
Set $\theta=\eta/2>0$.  Uniformly on the closed right unit semicircle,
\begin{align}
 \|Q_u^\sigma J(w)Q_u^\sigma\|&\le Cu,
 &\|Q_u^\sigma J(w)Q_u^{-\sigma}\|&\le Cu^\theta,
 \label{eq:angular-J-bands}\\
 Q_u^\sigma N(w)Q_u^\sigma
 &=\sigma Q_u^\sigma(wJ_{u,1}+w^{-1}J_{u,-1})Q_u^\sigma+O(u^2),
 \label{eq:angular-N-leading}\\
 \|Q_u^\sigma N(w)Q_u^{-\sigma}\|&\le Cu^{1+\theta},
 &\|Q_uN(w)Q_u\|&\le Cu.
 \label{eq:angular-N-bands}
\end{align}
The estimates remain valid for unequal input and output cutoffs, using
their maximum for $u$.  With $m$ ordinary time rows on each side,
$N$ costs $O(m^{-1})$, its remainder after the leading-mode expression
costs $O(m^{-1-\theta})$, and a sign-opposite undoubled $J$ costs
$O(m^{-\theta})$.

Finally, if $S=P(i)$ and $K$ is the principal argument of $S^2$, then
\begin{equation}\label{eq:angular-phase-bound}
 S^2=e^{iK},\qquad |K|\le C H_0.
\end{equation}
\end{lemma}

\begin{proof}
Evaluate \eqref{eq:general-inner-factorization} at $1$.  On the band
$H_0\le u$ it gives
\[
 a+b+\frac12\sum_j
 \log\left(1+\frac{4\Re\zeta_j}{|1-\zeta_j|^2}\right)\le u.
\]
The cone condition implies the zero budget
\begin{equation}\label{eq:angular-zero-budget}
 a+b+\sum_j\min(|\zeta_j|,|\zeta_j|^{-1})\le Cu.
\end{equation}
Indeed $\Re\zeta_j\ge c_0|\zeta_j|$ and
$|1-\zeta_j|\le1+|\zeta_j|$.  Each zero is therefore either of
radius at most $Cu$ or of reciprocal radius at most $Cu$.
Comparing each factor at $w$ and at $1$ shows, on
$Au<|w|<a/u$, $\Re w\ge0$, that
\[
 \left|\log\frac{f(w)}{f(1)}\right|
 \le Cu(|w|+|w|^{-1}).
\]
Choose $A$ large and $a$ small once and for all.  The function and its
inverse are then bounded on this half-annulus.  More precisely,
\begin{equation}\label{eq:angular-band-sign}
 P(w)Q_u^\sigma=\sigma Q_u^\sigma
       +O\bigl(u(|w|+|w|^{-1})\bigr),
\end{equation}
and the same estimate holds for the inverse restricted to the band.

The endpoint formulas give
\[
 J(i)=SZ_0,\qquad J(-i)=-Z_0S^{-1},\qquad
 J(-i)=-P(i)^{-1}J(i)P(i)^{-1}.
\]
The other endpoint gives the analogous identity on the other ray.
Dilate these uncompressed identities using
\eqref{eq:angular-covariance}, and only then compress.  This is
important because $Q_u$ need not commute with the dilation group.
The formula
\begin{equation}\label{eq:angular-current-gluing}
 J_u(-w)=-P(w)^{-1}J_u(w)P(w)^{-1},\qquad
 J_u(w)=Q_uJ(w)Q_u,
\end{equation}
defines the left half of the annulus; both inverses here and below are
only on $Q_u\mathcal H$.  Boundedness and strong boundary convergence
make the two definitions agree on both imaginary rays.  Scalar Morera
applied to every matrix coefficient gives holomorphic gluing.  Thus
$\|J_u(w)\|\le C$ on the full annulus.  Cauchy's formula on circles
near its inner and outer boundaries proves
\eqref{eq:angular-Laurent} and its coefficient bounds.  In particular,
uniformly on $|w|=1$,
\begin{equation}\label{eq:angular-odd-expansion}
 \tfrac12(J_u(w)-J_u(-w))
 =wJ_{u,1}+w^{-1}J_{u,-1}+O(u^3)=O(u).
\end{equation}
The nonconstant even modes are $O(u^2)$, but this observation alone
does not control the constant coefficient.

For $X_{\sigma\tau}(w)=Q_u^\sigma J(w)Q_u^\tau$,
\eqref{eq:angular-current-gluing} and
\eqref{eq:angular-band-sign} give
\[
 X_{\sigma\tau}(-w)=-\sigma\tau X_{\sigma\tau}(w)
                 +O(u)\|X_{\sigma\tau}(w)\|.
\]
On an equal-sign block its even part is therefore bounded by
$Cu\|X_{\sigma\sigma}(w)\|$.  Its odd part is $O(u)$ by
\eqref{eq:angular-odd-expansion}.  Absorbing gives
$\|X_{\sigma\sigma}\|\le Cu$ and then an $O(u^2)$ bound for its
whole even part, including the constant coefficient.  This proves the
first estimate in \eqref{eq:angular-J-bands}.

For opposite signs use an additional RTT comparison at $r=\sqrt u$.
The normalized intertwiner for arguments $r,w$ is
$R_\infty+O(r^\eta)$ with bounded inverse, uniformly for $w$ on the
closed right unit semicircle; $R_\infty$ is diagonal and invertible.
It commutes with the joint mark $I\otimes Z$.  Conjugating this mark
and expanding the $O(r^\eta)$ error in tensor matrix units, each
producing a bounded marked-row product, proves
\[
 \|[P(r),J(w)]\|\le Cr^\eta.
\]
Strong approximation includes the endpoints.  On $Q_u^\sigma$,
$P(r)=\sigma I+O(\sqrt u)$ by
\eqref{eq:angular-band-sign}.  Compression of the commutator yields
\[
 (2-C\sqrt u)\|X_{\sigma,-\sigma}(w)\|
 \le Cu^{\eta/2},
\]
which proves the second estimate in
\eqref{eq:angular-J-bands}.  Inserting
$P(w)=\sigma I+O(u)$ in \eqref{eq:angular-N-definition} now proves
\eqref{eq:angular-N-leading}--\eqref{eq:angular-N-bands}: equal signs
add, whereas opposite signs cancel.  In the latter case the odd part
of $X$ is itself $O(u^{1+\theta})$ by the gluing identity.  Thus the
leading expression can equivalently be written
\[
 \tfrac12\{\sgn(P(1))Q_u, wJ_{u,1}+w^{-1}J_{u,-1}\}
       +O(u^{1+\theta})
\]
on the whole band.  Higher modes need not vanish at a fixed cutoff;
the displayed error is what removes them after rescaling.

For unequal bands embed both in their larger band.  To obtain the
damped estimates, decompose each side into $H_0\le m^{-1}$, dyadic
intervals above $m^{-1}$ up to $u_0$, and the remaining spectrum.
On a band with lower endpoint $t$, $m$ physical rows contribute
$e^{-mt}$.  For any $\beta>0$ the double sum
\[
 \sum_{t,t'} C\max(t,t',m^{-1})^\beta e^{-cm(t+t')}
 \le C_\beta m^{-\beta}
\]
is finite with the asserted bound; unrestricted bounded rows handle
the exponentially damped remainder bands.  The same proof works when
the damping rows range in a fixed compact strict-angle interval, since
Schwarz--Pick gives a uniform comparison of their modulus loss with
$H_0$.

On a sufficiently small band, \eqref{eq:angular-band-sign} at $w=i$
gives $|\arg f(i)^2|\le C H_0(f)$, by applying the estimate with
cutoffs decreasing to the fiber value.  On its complement the principal
argument is at most $\pi$, which is also bounded by a fixed multiple
of $H_0$.  Spectral calculus proves
\eqref{eq:angular-phase-bound}.
\end{proof}

\subsection{The endpoint Hamiltonian and the sum rule}

The contraction $P(1)$ measures time damping.  We also need the generator
obtained by differentiating the unitary endpoint, because it is local
on spin states.  These are different operators: we keep the notation
$H_0$ for the former and $E$ for the latter.

\begin{lemma}[The local endpoint generator]\label{lem:endpoint-Hamiltonian}
The nonnegative self-adjoint spectral operator
\begin{equation}\label{eq:angular-endpoint-energy}
 E=S^{-1}\partial_zP(z)\big|_{z=v},\qquad
 E(f)=-\left.\partial_s\arg f(ie^s)\right|_{s=0},
\end{equation}
agrees on local-spin vectors with the vacuum-subtracted nearest-neighbor
Hamiltonian whose link is
\begin{equation}\label{eq:angular-local-link}
 h_{j,j+1}=-\eta\,\operatorname{swap}_{j,j+1}R'_{j,j+1}(0)
 =\begin{pmatrix}
 \eta\cot\lambda&0&0&0\\
 0&0&-\eta/\sin\lambda&0\\
 0&-\eta/\sin\lambda&0&0\\
 0&0&0&\eta\cot\lambda
 \end{pmatrix}_{j,j+1}.
\end{equation}
In particular $[h_{j,j+1},Z_j+Z_{j+1}]=0$ and $E\Omega=0$.
The closure of the local-spin subspace reduces $E$, and its real-time
evolution is the local Hamiltonian evolution.  For all sufficiently
small $u$,
\begin{equation}\label{eq:angular-E-H-comparison}
 \bbone_{\{E\le u\}}\le\bbone_{\{H_0\le Cu\}}.
\end{equation}
With $Z_1=S^{-1}Z_0S$, the exact adjacent-pair identity is
\begin{equation}\label{eq:angular-adjacent-current}
 N(i)=\tfrac12 S^2(Z_0+Z_1).
\end{equation}
\end{lemma}

\begin{proof}
Differentiate the inner factors along the imaginary boundary.  The
singular factors in \eqref{eq:general-inner-factorization} contribute
$a+b$, and a zero $\zeta$ contributes
\begin{equation}\label{eq:angular-zero-energy}
 \frac{2\Re\zeta}{|i-\zeta|^2}.
\end{equation}
All contributions are nonnegative.  Scalar continuation makes the
derivative finite almost everywhere and justifies the sum, or one may
first differentiate finite partial products and use positivity.
Since $\partial_z=iw\partial_w$, this is exactly the spectral
derivative in \eqref{eq:angular-endpoint-energy}.  The cone condition
bounds \eqref{eq:angular-zero-energy} below by a constant times
$\min(|\zeta|,|\zeta|^{-1})$.  Thus small $E$ implies the small/large
zero budget in \eqref{eq:angular-zero-budget}.  On that budget the
negative logarithm of each factor at $1$ is bounded above by the same
quantity.  This proves \eqref{eq:angular-E-H-comparison}.

At finite length, differentiate one factor in the swap row and divide
by its shift.  The inserted link is
$-\eta\operatorname{swap}R'(0)$, proving
\eqref{eq:angular-local-link}.  Dividing by the vacuum eigenvalue
subtracts its vacuum energy.  For a fixed local matrix $O$, the finite
action on $O\Omega_L$ is therefore
\[
 \left[\sum_jh_{j,j+1},O\right]\Omega_L,
\]
where only links meeting the support of $O$ occur.  Express
$O\Omega_L$ by endpoint rows, translating the support by a fixed shift
if necessary.  The finite division estimate and the RTT continuation
used in Lemma~\ref{lem:scalar-continuation} give uniform vector Taylor
bounds on a common endpoint disc for this fixed word.  Cauchy's formula
therefore passes its derivative, in norm and in all Gram pairings, to
the spectral derivative.  Consequently
\begin{equation}\label{eq:angular-local-commutator}
 E(O\Omega)=\left[\sum_jh_{j,j+1},O\right]\Omega.
\end{equation}
Iterating is legitimate because the commutator remains local.
If $O$ is supported on an interval of length $\ell$, each new
commutator enlarges the support by at most two.  Its norm is bounded
by $2\|h\|$ times the number of interacting links.  Hence
\[
 \|E^kO\Omega\|\le C_O C_O^k k!.
\]
Local vectors are analytic for the symmetric operator specified by
\eqref{eq:angular-local-commutator}.  Analytic-vector uniqueness makes
its closure self-adjoint on their closed span.  Since $E$ extends this
operator, the span reduces $E$ and the two self-adjoint operators agree
there.  This also identifies real-time evolution.  Finally
$J(i)=SZ_0$ and $P(i)=S$ in
\eqref{eq:angular-N-definition} give
\eqref{eq:angular-adjacent-current}.
\end{proof}

\begin{lemma}[Energy and winding sum rule]\label{lem:sum-rule}
Let $p$ denote the phase of the unit shift $S=P(i)$, modulo $2\pi$, and
let $\operatorname{sym}$ average a measure at $p$ and $-p$.  Set
\[
 D_0=-\langle\Omega,[Z_0,[Z_0,h_{0,1}]]\Omega\rangle.
\]
Away from $p=0$, the projection of the current measure obeys
\begin{equation}\label{eq:angular-sum-rule}
 2\operatorname{sym}\bigl((p)_*(E\mu)\bigr)
       =D_0\frac{\dd p}{2\pi}.
\end{equation}
On each orbit of degree $2n_F$, the phase winds monotonically through
$n_F$ complete turns.  In particular
\begin{equation}\label{eq:angular-degree-integral}
 \int n_F\dd\nu(F)<\infty.
\end{equation}
\end{lemma}

\begin{proof}
Since $Z_0\Omega=S^{-1}J(i)\Omega$, the spectral measure of the local
density vector is
\begin{equation}\label{eq:angular-density-measure}
 \|Z_0\Omega\|^2(\dd f)=|1-f(i)|^2\dd\mu(f).
\end{equation}
Fourier transform the vacuum double commutators of the translates
$Z_j$.  By $E\Omega=0$, their transform is twice the symmetrized
energy-weighted spectral measure of $Z_0\Omega$.  On the other hand
\eqref{eq:angular-local-commutator} and conservation of
$Z_0+Z_1$ leave only the three coefficients at separations $0,1,-1$,
namely $2D_0,-D_0,-D_0$.  Their transform is
$|1-e^{ip}|^2D_0$.  Thus
\[
 2\operatorname{sym}\bigl((p)_*(E|1-f(i)|^2\mu)\bigr)
 =|1-e^{ip}|^2D_0\frac{\dd p}{2\pi}.
\]
Testing on compact subsets of the circle without zero permits division
by the squared factor and proves \eqref{eq:angular-sum-rule}.

On a nonconstant finite-product orbit, direct differentiation gives
$-\dd p/\dd s=E(F(e^s\cdot))>0$.  The boundary phase change on the
positive imaginary ray is one half of the full boundary change: this
follows either by pairing conjugate zeros in
\eqref{eq:finite-inner-symmetry} or by integrating
\eqref{eq:angular-zero-energy} with its dilated argument.  It is
$2\pi n_F$.  At each nonzero phase there are consequently $n_F$
crossings, counted with multiplicity, and none has zero derivative.
The change of variables from $s$ to phase gives
\[
 (p)_*(E\,\dd s\,\nu)=\left(\int n_F\dd\nu(F)\right)\dd p
\]
away from zero, first with truncated orbit measures and then by monotone
convergence.  Equation~\eqref{eq:angular-sum-rule} implies finiteness;
in these conventions it also gives
$4\pi\int n_F\dd\nu=D_0$.  The endpoints of an orbit and its finitely
many zero-phase crossings have zero $\dd s$ measure, so no missing
phase atom is involved in this computation.
\end{proof}

\subsection{The height variogram}

The measure constructed above already determines height covariances.
The following bound will also supply the logarithmic majorants needed
for nonsmooth tests and regularity.

\begin{proposition}[Exact variogram and logarithmic growth]
\label{prop:variogram}
After an axial symmetry, any nonzero lattice displacement may be
written $(k,j)$ with $k\ge1$ and $|j|\le k$.  Its height variance is
\begin{equation}\label{eq:angular-variogram}
 \begin{aligned}
 \E_c\bigl(h(k,j)-h(0,0)\bigr)^2
 &=2\Re\int\bigl(1-f(1)^k f(\pm i)^{|j|}\bigr)\dd\mu(f)\\
 &=4\kappa\log k+O(1),
 \end{aligned}
\end{equation}
where the sign of $i$ is chosen according to the vertical direction.
The error is uniform over these displacements.  In particular the
variance is bounded by $C(1+\log(1+|k|+|j|))$.
\end{proposition}

\begin{proof}
First consider a directed path of physical or endpoint steps, and write
their spectral multipliers as $t_1,\ldots,t_m$.  Reflection changes
the sign of a current.  Equations~\eqref{eq:current-cocycle} and
\eqref{eq:angular-endpoints} therefore express the covariance of two
distinct increments, in their path order, as the negative real part of
the integral of
\[
 (1-t_a)t_{a+1}\cdots t_{b-1}(1-t_b),\qquad a<b,
\]
with the conjugate term supplying the real covariance.  A single
increment has variance one.  At an endpoint this equals
$2\Re\int(1-f(i))\dd\mu$, because $|f(i)|=1$ and
\eqref{eq:angular-density-measure} integrates to
$\|Z_0\Omega\|^2=1$.  The same identity holds at the physical step
$w=1$: on each orbit move the ray from $i\R_+$ to $\R_+$ by Cauchy's
theorem for $(1-F(w))\dd w/w$.  Its small and large circular arcs
vanish, since $F(0)=F(\infty)=1$.  The integrated identity is
justified by the bounds below.

The elementary telescoping identity
\[
 1-\prod_{a=1}^m t_a
 =\sum_{a=1}^m(1-t_a)
  -\sum_{a<b}(1-t_a)t_{a+1}\cdots t_{b-1}(1-t_b)
\]
therefore gives the first equality of
\eqref{eq:angular-variogram}, taking $k$ horizontal and $|j|$
vertical steps.  It is equally valid by truncating the orbit measure
first and then passing to the absolutely convergent integrals below.

Here are bounds that both justify this passage and give the asymptotic.
On an orbit choose one zero radius as origin in logarithmic coordinates
and let $\mathcal L$ be the range of the log radii.  Beyond either
extreme by distance $d$, the factor formulas and the endpoint
normalizations give
\begin{equation}\label{eq:angular-orbit-tail}
 |1-F(e^s)|+|1-F(\pm i e^s)|\le Cn_F e^{-d},
\end{equation}
or the trivial constant bound when it is smaller.  The factors on the
imaginary ray have modulus one; those on the positive ray have modulus
at most one.  Hence the integrand in
\eqref{eq:angular-variogram} is absolutely integrable on each orbit.
Outside the zero range enlarged by $\log k+\log(1+n_F)$, the inequality
$|1-a^kb^{|j|}|\le k|1-a|+|j||1-b|$, with $|j|\le k$, and
\eqref{eq:angular-orbit-tail} give a bounded integral.  The remaining
extra interval lengths cost at most
$C(\mathcal L+\log(1+n_F)+1)$.

For the main interval, use any single zero.  Its cone bound gives
\[
 |F(e^s)|^k\le
 \exp\{-c k e^{-|s-s_0|}\}
 \quad\text{for } |s-s_0|\le\log k,
\]
where $s_0$ is its logarithmic radius (with the sign dictated by the
chosen orbit coordinate).  The integral of the right side over this
interval is bounded uniformly in $k$.  Multiplication by
$F(\pm i e^s)^{|j|}$ does not change that bound.  Consequently the
integral of the real part of $1-F(e^s)^kF(\pm i e^s)^{|j|}$ is
\[
 2\log k+O\bigl(1+\mathcal L+\log(1+n_F)\bigr).
\]
The remainder is integrable against $\nu$ by
\eqref{eq:orbit-width-tail}, \eqref{eq:angular-degree-integral}, and
$\kappa<\infty$.  Multiplication by the factor $2$ in the variogram
proves \eqref{eq:angular-variogram}.  The same majorants justify the
orbitwise contour move and the telescoping passage above.  Axial
symmetry, already proved for the physical law, supplies the final
uniform displacement bound.
\end{proof}

We have thus obtained a physical representation with a finite measure
$\nu$ of dilation orbits, an as yet unevaluated constant $\kappa$, and
uniform low-energy current estimates.  Section~\ref{sec:normalization}
determines $\kappa$.  Section~\ref{sec:chiral} uses the same estimates,
including their higher-mode errors, to identify all limiting current
correlations.
